\documentclass[serif]{beamer}
\input{sns_beamer_preamble.tex}

\begin{document}

\begin{frame}{중괄호로만 표현하기}
    위수가 $2$인 순환군 $C_2$를 중괄호와 쉼표만으로 표현해 봅시다.
    \[
        \begin{aligned}
            C_2&=(\{0,1\},+)\\
            &\phantom{=}\quad\quad\quad\quad\quad\quad\quad\quad\quad\quad\quad\quad\quad\quad\quad\quad\quad\quad\quad(\text{집합과 이항 연산의 순서쌍})\\
            &=\{ \{ \{0,1\} \}, \{ \{0,1\}, +\} \}\\
            &\phantom{=}\quad\quad\quad\quad\quad\quad\quad\quad\quad\quad\quad\quad(\text{Kuratowski의 정의: }(a,b)=\{ \{a\}, \{a,b\} \})\\
            &=\{ \{ \{0,1\} \}, \{ \{0,1\}, 
                \{
                    ((0,0),0), ((0,1),1), ((1,0),1), ((1,1),0)
                \}
            \} \}\\
            &\phantom{=}\quad\quad\quad\quad\quad\quad\quad\quad\quad(\text{이항 관계로 다시 쓴 이항 연산 }+:\{0,1\}^2\to\{0,1\})\\
            &=\{ \{ \{0,1\} \}, \{ \{0,1\}, 
                \{
                    \{ \{(0,0)\}, \{ (0,0), 0 \} \},
                    \{ \{(0,1)\}, \{ (0,1), 1 \} \},\\
            &\phantom{=}\quad\quad\quad\quad\quad\quad\quad\quad\quad
                    \{ \{(1,0)\}, \{ (1,0), 1 \} \},
                    \{ \{(1,1)\}, \{ (1,1), 0 \} \}
                \}
            \} \}\\
            &\phantom{=}\quad\quad\quad\quad\quad\quad\quad\quad\quad\quad\quad\quad\quad\quad\quad\quad\quad\quad\quad\quad\quad(\text{Kuratowski의 정의})\\
            &=\{ \{ \{0,1\} \}, \{ \{0,1\}, 
                \{
                    \{ \{ \{0\} \}, \{ \{0\}, 0 \} \},
                    \{ \{ \{ \{0\}, \{0,1\} \} \}, \{ \{ \{0\}, \{0,1\} \}, 1 \} \},\\
            &\phantom{=}\quad\quad\quad\quad\quad\quad\quad
                    \{ \{ \{ \{1\}, \{1,0\} \} \}, \{ \{ \{1\}, \{1,0\} \}, 1 \} \},
                    \{ \{ \{1\} \}, \{ \{1\}, 0 \} \}
                \}
            \} \}\\
            &\phantom{=}\quad\quad\quad\quad\quad\quad\quad\quad\quad\quad\quad\quad\quad\quad\quad\quad\quad\quad\quad\quad\quad(\text{Kuratowski의 정의})\\
            &=\{\{\{\{\},\{\{\}\}\}\},\{\{\{\},\{\{\}\}\},\{\{\{\{\{\{\}\}\}\},\{\{\{\{\}\}\},\\
            &\phantom{=}\quad\{\}\}\},\{\{\{\{\{\}\},\{\{\},\{\{\}\}\}\}\},\{\{\{\{\}\},\{\{\},\{\{\}\}\}\},\\
            &\phantom{=}\quad\{\{\}\}\}\},\{\{\{\{\{\{\}\}\},\{\{\{\}\},\{\}\}\}\},\{\{\{\{\{\}\}\},\{\{\{\}\},\{\}\}\},\\
            &\phantom{=}\quad\{\{\}\}\}\},\{\{\{\{\{\{\}\}\}\}\},\{\{\{\{\{\}\}\}\},\{\}\}\}\}\}\}\\
            &\phantom{=}\quad\quad\quad\quad\quad\quad\quad\quad(\text{von Neumann의 정의: }0=\{\},1=\{0\},2=\{0,1\},\cdots)
        \end{aligned}
    \]
    이렇게 148개의 중괄호로 $C_2$를 나타낼 수 있습니다. ($\{$과 $\}$를 각각 셉니다.)
\end{frame}

\begin{frame}{}
    \begin{hint}[문제]
        같은 규칙으로 위수가 $3$인 순환군
        \[
            C_3=(\{0,1,2\},+)
        \]
        는 몇 개의 중괄호로 나타낼 수 있을까?
    \end{hint}
\end{frame}
\end{document}